This paper argues that optimization is not a law of the universe, and that morality and the manner of being of intimate relation are not, in their structure, optimization problems. It proceeds in two movements that it is careful to keep apart. The first is a formal proof, carrying no value judgment, that intimate relation furnishes no well-defined optimization problem: the value at stake is in part non-structurable, since the drive, which any complete objective must include, resists representation as a term; it is relational, generated within the relation and existing only there, so that no external standpoint could fix the objective; it is incommensurable, so that to price certain goods is to betray rather than mismeasure them; and the subject is constituted within the relation rather than given before it, so that the optimizer the model presupposes is a product of the process it is meant to stand outside of. Each failure strikes a distinct precondition, and any one leaves the operation of maximization undefined.

The second movement takes that result as given and criticizes the imposition of the frame regardless. To optimize where no optimization problem exists converts a value that should circulate into one that can be settled and drawn off, in the manner the theory of generative justice names as extractive computation; it is a symptom of the crisis of the symbolic and of a life-world habituated to the logic of capital, under which a local, historical operation is projected onto the cosmos as its law; and it mistakes the kind of thing morality is, since an optimizing practice, even a successful one, may fail as an ethics, the optimal and the good being orthogonal. The paper then recovers a rationality that responds rather than maximizes, described through response, attunement, and the sustaining of generation, and it closes in polyphony, refusing to name a new master objective, since to do so would reinstate the frame. That refusal is the paper's positive claim, enacted in its form.
